\section{Conclusion}
\label{sec:conclusion}

We asked when and why the cosine between an edit's key and an unrelated
probe's key predicts locate-then-edit collateral damage, and answered with a
mechanism: for rank-one ROME updates the perturbation factorizes in closed
form into edit strength, probe norm, and key-cosine -- the exact rank-one
reduction of first-order gradient influence, a zero-cost rank estimator that
is nonetheless not a faithful rank-surrogate of the true influence it reduces.
On the confound-clean within-probe metric the law traces an inverted-U in
depth, is overtaken by norm-growth at the deepest layer studied, and its sign
tracks the damage regime across scale with attenuated magnitude at the largest
size tested; the editable band itself is architecture-conditioned, from wide
plateau to shallow cliff.

That conditionality is part of the contribution: the law vanishes for
full-rank fine-tuning and multi-layer-spread editors and goes null or inverts
off the Llama family, yet its causal reality stands -- null-space projection
removes the geometry-predicted damage in proportion to key-cosine at every
layer studied, erases the one architecture's sign inversion, and generalizes
in form to an instruction-tuned twin and a second architecture (both
three-seed) and at 20B (single-seed), while the signed coupling attenuates
at 8B, scoping the
strong positive-signed law to small-to-mid scale. The mechanism extends to
refusal-deletion edits and replicates across datasets, including multi-hop
MQuAKE under an overlap-robustness audit -- complementing the concurrent
empirical predictor of Section~\ref{sec:related} with the mechanism, its
conditions, and the causal test closing the loop.

Practically, key-cosine geometry is a zero-cost diagnostic for anticipating
and, via null-space projection, mitigating collateral damage -- within the
scope delineated here: a specific editor family, architecture family, and depth
range short of the regime transition. Extending the causal test to larger
scales, more editors, and more architectures is the natural next step
(Section~\ref{sec:discussion}).
